Set
\[
\varepsilon_0=\varepsilon_1=\alpha/2.
\]
By \(\texttt{lem:armwise-finite-rank-staircase}\),
\[
k_0=k_1=k,
\]
and \(\texttt{def:coverage-frontier}\) gives
\[
\Gamma_{M,\alpha,\boldsymbol\varepsilon}(w)=\Gamma_{M,\alpha}(w).
\tag{1}
\]
Therefore \(\texttt{thm:asymmetric-sharp-orbit-coverage-frontier}\) applies on this diagonal. Its total-preorder argument proves exactness over the declared bounded orbit-constant score class and attainment by \(F_{m_0}F_{m_1}\) preorder pairs, with at most \(m_a\) score levels in arm \(a\). It also preserves the unchanged coherent \(P_T^C\), full joint \(R_w\), local weights \(\omega_C\), common target, and common selector.

Suppose \(k\le M\). Put
\[
\bar k=\min\{k,M-k+1\},\qquad
\mathsf C_{M,k,k}
 =\sum_{i=0}^{M-1}\min\{\bar k,i+1\}^2.
\tag{2}
\]
Under the substitution \(k_0=k_1=k\), the asymmetric theorem's exact-key partition is
\[
\Pi_\star
 =\Pi_{\mathrm{roll}}\mathbin{\dot\cup}\Pi_{\mathrm{MS}}\mathbin{\dot\cup}\Pi_+,
\tag{3}
\]
where every key is deterministically assigned to a least certified charge among
\[
\mathsf C_{M,k,k},\qquad
M+\mathsf T_{M,\bar k,\bar k},\qquad
\mathsf C_+(\pi).
\]
Thus its total key-evaluation charge specializes to
\[
\begin{aligned}
\mathsf B_{\mathrm{sym}}={}&
M\mathbf 1\{\Pi_{\mathrm{MS}}\ne\varnothing\}
+\sum_{\pi\in\Pi_{\mathrm{roll}}}\mathsf C_{M,k,k}\\
&+\sum_{\pi\in\Pi_{\mathrm{MS}}}
 \left(M+\mathsf T_{M,\bar k,\bar k}\right)
+\sum_{\pi\in\Pi_+}\mathsf C_+(\pi),
\end{aligned}
\tag{4}
\]
and its selected-branch workspace specializes to
\[
\mathsf W_{\mathrm{sym}}
=\max\!\left(
\{0\}
\cup\{\bar k^2:\Pi_{\mathrm{roll}}\ne\varnothing\}
\cup\{M:\Pi_{\mathrm{MS}}\ne\varnothing\}
\cup\left\{\min\{t_0^{\sigma(\pi)},t_1^{\sigma(\pi)}\}:\pi\in\Pi_+\right\}
\right).
\tag{5}
\]
Substituting (4)--(5) into the asymmetric instance-adaptive bound yields
\[
O\!\left(
|\mathcal V|+F_{m_0}F_{m_1}(m_0m_1+L)
+\sum_{C\in\mathfrak C}\!\left[L_C\log(1+L_C)+L_C+U_C\right]
+\left\{\sum_{C\in\mathfrak C}U_C\right\}\log(1+G)
+\mathsf B_{\mathrm{sym}}
\right).
\tag{6}
\]
Replacing \(L_C,U_C\) by \(L_{\max},U_{\max}\) gives the stated uniform bound, and the memory bound becomes
\[
O(m_0m_1+L+U_{\max}+G+\mathsf W_{\mathrm{sym}}).
\tag{7}
\]
The first term in (4) is present exactly once if the multinomial branch is nonempty, because factorial preprocessing is shared. Every exact key remains a function of the unchanged \(R_w\), and every local weight remains a function of the unchanged coherent \(P_T^C\), so (3) changes computation but not the common-target or cross-arm probability. At \(k=M\), one has \(\bar k=1\) and
\[
\mathsf C_{M,M,M}=M.
\]
Hence rolling costs \(\Theta(M)\) per key rather than the coarse \(\Theta(M^3)\), with a potentially smaller positive-corner coefficient charge.

If \(k=M+1\), both armwise rank events are certain by \(\texttt{prop:rank-boundary}\), so
\[
\Gamma_{M,\alpha}(w)=[(1-\alpha)-1]_+=0.
\tag{8}
\]
The admissible four-pair construction in the asymmetric theorem already lies on the diagonal \(\varepsilon_0=\varepsilon_1=\alpha/2\). It gives the four-level objective
\[
\frac{24957}{160000},
\]
the binary upper bound
\[
\frac{759109}{5000000},
\]
and the strict gap
\[
\frac{83189}{20000000}.
\]
Thus binary attainment is false in the symmetric problem as well. Finally, the asymmetric disjoint-support construction specializes through (1): whenever one arm has disjoint target and calibration orbit supports and \(k\le M\), its rank event fails surely, so the defining deficit equals \(1-\alpha\).